\section{Realizing the scalar conditions}\label{sec:realization}

Write \(r(v)=a+T(v,\cdot)\). A hole requires the opposite vertex's
channel functional to equal \(r(v)\) on the whole profile space.
The lists determine its values on the selected atoms. The entries
already fixed determine its values on the pins. We first find the
conditions under which these prescriptions agree, then extend them
to the whole space.

\subsection{Lists and their records}

Let \(A=(A_1,A_2)\) and \(B=(B_1,B_2)\) be two units. For each pair
\(i,j\in\{1,2\}\), choose a nonempty list of at most 7 atoms at
\(A_i\), and a matching list at \(B_j\). Matched atoms have the same
tag and diagonal bit. Let \(v^A_{ij},v^B_{ji}\) be their sums.
There are at most 28 atoms on either unit, and at most 14 at one endpoint.
Their selector and base directions satisfy the independence conditions
of Lemma~\ref{lem:protected-atoms}.

The key of \(A(l,s,z)\) at an orientation \(o\) is
\[
 \bigl(P_{e_{l-1}}v(s,z),P_{e_l}v(s,z)\bigr).
\]
Equality of matched keys implies
\(U_{A_i}v^A_{ij}=U_{B_j}v^B_{ji}\).
Define the old internal functional at endpoint \(i\) of a unit by
\[
 p_i=u_{A_{3-i}}U_{A_i}.
\]
Here \(p_i\) evaluates an atom using the channels of the other
endpoint of the \emph{same} unit. It is one of the bits recorded before
two units are compared.

The pins of the two leaves determine a bounded \emph{table} of cross
pairings. It records primal pin pairings with the opposite channels, and
the entries needed to keep all stored images fixed. The flags on a unit
say that its entries agree with that table. Because primal projections
are stored pins, these entries determine the functional
\(u_{B_j}U_{A_i}\) on the effective space \(C_i^A\).
Call that functional \(t^A_{ij}\). Define \(t^B_{ji}\) reciprocally.

The numerical conditions on the lists are
\begin{align}
 \sum_{a\text{ in list }ij}q_a&=1,
 &a(v^A_{ij})+a(v^B_{ji})&=1,\label{eq:recipe-scalars}\\
 T(v^A_{ij},\cdot)|_{C_i^A}&=(t^A_{ij}+a)|_{C_i^A},
 &T(v^B_{ji},\cdot)|_{C_j^B}&=(t^B_{ji}+a)|_{C_j^B}.
 \label{eq:recipe-gradients}
\end{align}
At each opposite endpoint, the two lists must also have equal sums of
\(p+a\). For example, on unit \(A\),
\begin{equation}\label{eq:recipe-balance}
 (p_1+a)(v^A_{1j})=(p_2+a)(v^A_{2j})\quad(j=1,2).
\end{equation}
There is the reciprocal condition on \(B\).

% lax begin Lax793645.ScalarPrescriptions
\begin{lemma}\label{lem:small-tables}
If \eqref{eq:recipe-scalars}--\eqref{eq:recipe-balance} hold, the binary
values between selected atoms can be prescribed so that
\(r(v^A_{ij})\) has the correct value on every selected atom at \(A_i\),
and likewise at the other endpoints. These prescriptions use only the
reduced mixer family.
\end{lemma}
% lax end
% lax begin Lax793645Proofs.Manuscript.ScalarPrescriptions.choose_reduced_script
\begin{proof}
Consider the matrix of pairings between all selected atoms at one endpoint.
Its off-diagonal entries may be chosen subject to symmetry. Its diagonal is the role bit,
since \(T(x,x)=a(x)\). Prescribing the sums along the two outgoing lists
amounts to prescribing the products of this matrix with the two list
indicator vectors. The diagonal identity imposes the self-pairing
conditions in \eqref{eq:recipe-scalars}. Symmetry imposes equality of the
two cross sums, which is \eqref{eq:recipe-balance} on the opposite unit.
The two list indicators have disjoint nonempty supports, so they extend
to a basis. In this basis fill their prescribed rows and columns.
The cross sums give the same value at the entry they share. On the
remaining basis vectors set the diagonal to the prescribed linear
functional and choose the other entries symmetrically. This works
because in characteristic 2 the diagonal of a symmetric bilinear
form is a linear functional.

It remains to realize each chosen scalar pairing using \(T^1\).
For tags \(l,l'\), one of the two cycle edges incident with \(l\) has
index \(c\ne l'\). Equation \eqref{eq:matching-atom} shows that its
indexed term couples these atoms. Use that term for the prescribed value
and set the other products to zero. Distinct ordered atom pairs have
separate binary prescriptions. The explicit binary mixing estimate applies
to the independent indexed tests used in this prescription.
\end{proof}
% lax end

\subsection{Completing the channel coordinates}

Let \(H=\F^h\oplus\F^h\), with \(q_H(x,y)=x\cdot y\).
The subspace \(F\le B'\) consists of directions whose images
\(g_0:F\to H\) are already fixed. The subspace \(C\le H\) plays
the complementary role in the target. The pairings against \(C\)
are prescribed for \emph{every} input by a linear map
\(\ell:B'\to C^*\).
We seek a map \(g:B'\to H\) with prescribed quadratic value \(Q\),
agreeing with \(g_0\), and satisfying
\(\beta_H(gb,c)=\ell(b)(c)\).

% lax begin Lax793645.CoordinateCompletion
\begin{lemma}[Coordinate completion]\label{lem:relative-extension}
Such a map exists if
\begin{align*}
 Q(f)&=q_H(g_0f), &
 \ell(f)(c)&=\beta_H(g_0f,c),\\
 \beta_Q(f,b)&=\ell(b)(g_0f)
 &&\text{whenever }f\in F,\ g_0f\in C,
\end{align*}
and \(Q(b)=b\trans Mb\), with
\[
 h\ge\rank M+5\dim C+6\dim F+2.
\]
The map \(g_0\) may have a kernel.
\end{lemma}
% lax end
% lax begin Lax793645Proofs.Manuscript.CoordinateCompletion.coordinate_quadratic_extension
\begin{proof}
Write \(f=\dim F\), \(c=\dim C\), and \(\rho=\rank M\).
First prescribe the pairing of \(g_1b\) with \(C+g_0F\) by
\[
 \beta_H(g_1b,c_0+g_0f_0)=\ell(b)(c_0)+\beta_Q(b,f_0).
\]
There can be two ways to write the same vector as \(c_0+g_0 f_0\).
Their difference has \(c_0=g_0 f_0\in C\), and the last compatibility
condition makes the right side zero on every such relation.
Nondegeneracy of \(\beta_H\) makes \(H\to(C+g_0F)^*\) surjective.
Choose a linear section. On \(F\) its solution can differ from \(g_0\)
only by a vector orthogonal to \(C+g_0F\). Extend this difference by
zero on a complement of \(F\) and add it to the solution. The resulting
map \(g_1\) agrees with \(g_0\). The section contributes at most \(c+f\)
to its rank and the correction at most \(f\).

Put \(R=Q+q_H\circ g_1\). For \(f_0\in F\), the prescribed quadratic
value gives \(R(f_0)=0\), and the prescribed pairing with \(g_0 f_0\)
gives \(\beta_R(b,f_0)=0\). Thus \(R(b+f_0)=R(b)\). Project \(B'\)
onto a complement of \(F\) along \(F\). A matrix for \(q_H\circ g_1\) is
\((g_1^+)\trans g_1^-\), where \(g_1^+,g_1^-\) are the coordinate
halves of \(g_1\). Its rank is at most \(c+2f\). Adding \(M\) and
pulling back by the projection still represents \(R\), with rank at
most \(\rho+c+2f\).

Rank factorization therefore writes
\[
 R(b)=\sum_{j=1}^k a_j(b)b_j(b),\qquad k\le\rho+c+2f,
\]
where all the linear forms vanish on \(F\).

Put \(W=C+g_1B'\), of dimension \(d\le2c+2f\).
In the coordinate half \(\F^h\oplus0\), choose \(k\) vectors \(x_j\)
in \(W^\perp\), independent modulo \(W\). Orthogonality to \(W\) imposes at most \(d\) equations in this
\(h\)-dimensional half, and quotienting by its intersection with \(W\)
loses at most another \(d\) dimensions. Thus \(h\ge k+2d\) suffices.
Independence modulo \(W\) makes the following equations soluble. Choose
\(y_j\in W^\perp\) with \(\beta_H(x_i,y_j)=\delta_{ij}\).
Adding \(q_H(y_j)x_j\) makes \(q_H(y_j)=0\). Adding suitable preceding
\(x_i\)'s then makes all pairings between the \(y_j\)'s zero, without
changing those quadratic values or the prescribed pairings.
Thus
\[
 g(b)=g_1(b)+\sum_j\bigl(a_j(b)x_j+b_j(b)y_j\bigr)
\]
has quadratic value \(Q(b)\). The correction vanishes on \(F\) and is
orthogonal to \(C+g_1B'\). It preserves every linear constraint.
Finally \(k+2d\le\rho+5c+6f\), which proves the bound.
\end{proof}
% lax end

The kernel condition matters here. A law obeying the density bounds can
give selected keys zero channel image. For such an \(f\), every extension
has \(\beta_Q(f,b)=0\) for all \(b\). The last compatibility equation
imposes exactly this requirement.

\subsection{The obstruction calculation}

We describe the calculation jointly for the two endpoints of one unit.
At a component the domain is \(B'=B\oplus B\). Let
\(F_i=P_i+A_i\), where \(P_i\) is the primal pin space and \(A_i\)
is the selected-key span. Put \(F=F_1\oplus F_2\).
For one opposite endpoint, let \(g_0:F\to H\) be its already fixed
channel evaluation and put
\[
 K_0=\{f\in F:g_0f\in C\}.
\]
The table is called \emph{safe} if
\begin{equation}\label{eq:endpoint-injection}
 K_0\cap(P_1\oplus0)=0
 \quad\text{and}\quad K_0\cap(0\oplus P_2)=0.
\end{equation}
That is, its
channel map is injective on each endpoint's primal pins modulo the stored
channel space. No injectivity is required on their sum.

The homogeneous quadratic constraints in Lemma~\ref{lem:relative-extension}
have annihilator
\begin{equation}\label{eq:quadratic-annihilator}
 F^{(2)}+K_0\odot B',\qquad
 K_0\odot B'=\Span\{kb\trans+bk\trans:k\in K_0,b\in B'\}.
\end{equation}
Indeed evaluation on \(ff\trans\) reads the frozen quadratic value,
and evaluation on \(kb\trans+bk\trans\) reads a polar row.
Every matrix in this annihilator lies in
\(F\otimes B'+B'\otimes F\), so its rank is at most \(2\dim F\).

% lax begin Lax793645.FrozenConstraints
\begin{lemma}\label{lem:joint-obstruction}
An annihilating profile for the frozen quadratic constraints is a sum of
selected atoms and effective profiles on the primal pins.
\end{lemma}
% lax end
% lax begin Lax793645Proofs.Manuscript.FrozenConstraints.one_selected_obstructions
\begin{proof}
Fix a component at endpoint \(i\), and let \(H_i\) be the column
space of its symmetric matrix \(w_i\). Its rank is at most \(R\), by
\eqref{eq:quadratic-annihilator} and the choice of \(R\).
Lemma~\ref{lem:selector-confinement} bounds the reshaped rank of every
\(h\in H_i\) by \(R^2\). In particular, if \(h\in H_i\cap F_i\), write
\(h=p+a\) with \(p\in P_i\) and \(a\in A_i\). The vector \(a\) is a
sum of at most \(t_*\) simple tensors. Thus the reshaped rank of \(p\)
is at most \(R^2+t_*\), so \(p\in D_0\).

By symmetry \(w_i\in H_i\otimes H_i\). Its image vanishes when both
factors are taken modulo \(F_i\), by \eqref{eq:quadratic-annihilator}.
On \(H_i\) the quotient map has kernel \(H_i\cap F_i\). Its tensor
square therefore has kernel
\[
 (H_i\cap F_i)\otimes H_i+H_i\otimes(H_i\cap F_i).
\]
We just proved \(H_i\cap F_i\subseteq D_0+A_i\).
Lemma~\ref{lem:protected-atoms} now separates the selected diagonal
matrices from a remainder whose columns lie in \(L_i\).

The coefficients of a selected atom agree on its two incident components.
To check this, project the cut relation \(\sum_ex_e=0\) modulo the protected
spaces. The selected selector and base vectors are independent in these
quotients. A particular selected label appears on just its two incident
components. Its two coefficients must therefore sum to zero. We have
therefore subtracted genuine atoms, and the residual is still a profile.

It remains to prove that the residual is supported on pins. This requires
care with dependencies between endpoints. Write its component as
\[
 W=\operatorname{diag}(W_1,W_2)=M+T+T\trans,
 \qquad M\in F^{(2)},\quad\operatorname{col}T\subseteq K_0.
\]
If \(u\) annihilates \(F_i\), insert \(u\) in endpoint \(i\) and zero
in the other endpoint. The terms \(M\) and \(T\trans\) kill this vector.
Hence its image under \(W\) lies in \(K_0\). Block diagonality places
the image at endpoint \(i\), where its column space lies in \(L_i\).
Since \(K_0\le F\), its nonzero block lies in
\(F_i\cap L_i=P_i\). Safety \eqref{eq:endpoint-injection} forces it to
be zero. Symmetry now puts all columns of \(W_i\) in \(F_i\), and thus
in \(P_i\). A symmetric tensor supported there belongs to \(P_i^{(2)}\).
This is precisely an effective profile.
\end{proof}
% lax end

A kernel vector may still involve pins at both endpoints. The argument
above applies the injection condition only to a vector supported at one
endpoint, which is why the two separate conditions in
\eqref{eq:endpoint-injection} suffice.

The target functional has the correct values on the obstructions just
identified. Equation \eqref{eq:recipe-gradients} gives agreement on
effective profiles, and Lemma~\ref{lem:small-tables} gives agreement
on the selected atoms. Finite-dimensional duality
therefore gives a quadratic target satisfying all frozen restrictions
and all required profile evaluations.

\subsection{Keeping the target rank bounded}

A quadratic target exists, but arbitrary solutions could have rank growing
with \(n\). We need a solution of bounded matrix rank for coordinate completion.
The explicit functional \(r(v)\), for a list of at most 7 atoms whose
diagonal sum is one, has a bounded-rank matrix representative.
For the tester part, the ordinary coefficient at index \(t\) is the shared
tester value of the given atom-sum representation. At most 7 indices
contribute. Equation \eqref{eq:cycle-functionals} puts the shared terms
on the matchings \(M_t\). After collecting them, each component carries
at most one copy of the shared tester as well as at most 7 ordinary
testers. Together their rank on that component is at most \(8r_0\).
Including the \(T^1\) terms on both endpoints, we use the bound
\begin{equation}\label{eq:gradient-rank}
 r_{\rm grad}=2(15r_0+14Jg).
\end{equation}
Here the term \(14Jg\) bounds the contribution of the matrices in
\eqref{eq:reduced-mixer} at one endpoint, and \(15r_0\) bounds its tester
contribution.

Let \(f=\dim F\), \(c=\dim C\). A baseline channel map extending
\(g_0\) and realizing \(\ell\) can be chosen of rank at most \(f+c\).
Lift \(\ell\) through \(H\to C^*\) and extend the remaining correction
from \(F\). Its quadratic form satisfies every frozen constraint.
Project the matrix representing the prescribed functional away from \(F\). The difference from its
original matrix has rank at most \(2f\). Subtract the baseline and this
projected gradient. The residual has zero frozen values and zero polar
rows on \(K_0\), and its required profile functional has a representative
of rank bounded in terms of \(f,c\).

We reduce the rank of the residual by keeping just the base directions
on which its prescribed values depend.
In each base block choose a subspace containing the factors of all frozen
vectors and of the rows and columns in the bounded-rank target. Choose a
linear retraction that fixes the former and whose transpose fixes the
latter.

To construct it, let \(U\) be the space of vectors to fix and \(D\) the
space of linear forms to fix. Take a complement of \(U\cap\ker D\) in
\(\ker D=\bigcap_{\phi\in D}\ker\phi\) as the kernel of the retraction.
This kernel is disjoint from \(U\). Project along it onto a complement
containing \(U\). The rank is at most \(\dim U+\dim D\), and
\(\phi\) agrees with its pullback for every \(\phi\in D\).

Keep the constant
coordinate fixed and apply the same base map at every tag and endpoint.
This map preserves the allowed coordinate blocks, hence the moment spans
and the cut relations. Pulling the residual quadratic form back along it
preserves its profile values. Since it fixes \(F\), it preserves the zero
frozen values and the zero polar rows on \(K_0\).

The compressed residual has rank at most a fixed number \(d\), independent
of \(n,r_0,h\). Restoring the baseline and explicit part gives matrix rank
at most \(r_{\rm grad}+f+c+d\). It is enough to have
\begin{equation}\label{eq:channel-room}
 h\ge r_{\rm grad}+7f+6c+d+2.
\end{equation}
The coefficient of \(r_0\) on the right is 30. All other terms are fixed
before \(r_0\). Thus \(h=1000r_0\) suffices for large \(r_0\).

% lax begin Lax793645.FrameRealization
\begin{theorem}[Realization]\label{thm:gradient-realization}
For protected lists satisfying the scalar conditions, a consistent safe
table and matching keys admit a completion of the cross pairings that
realizes all four hole functionals. The completion preserves every frozen
row and column. It prescribes at most \(16gh\dim B\) further scalar entries.
\end{theorem}
% lax end
% lax begin Lax793645Proofs.Manuscript.FrameRealization.one_reciprocal_realization
\begin{proof}
Use the joint obstruction calculation and compression for each opposite
channel map, then apply Lemma~\ref{lem:relative-extension}. The reciprocal
problems occupy different blocks. One prescribes the pairings of the
first unit's primal columns against the other unit's channels. The
reciprocal problem prescribes the pairings with the two roles reversed. Their shared
frozen entries already agree with the table. They can therefore be completed
independently. The bound \(16gh\dim B\) counts pairings between primal and channel
columns in both directions over all endpoint pairs and components. Once their actual values equal the
completion, \eqref{eq:recipe-scalars} and the matching keys give the hole
relations for all four pairs.
\end{proof}
% lax end
